\begin{figure}[ht]
\centering
\begin{tikzpicture}[x=1.0cm, y=1.0cm, font=\small]

% Panel 1: input unit circle in R^2
\begin{scope}[shift={(0,0)}]
  \node[anchor=south, font=\small\bfseries] at (0, 2.0) {input: unit circle};
  \draw[->, thick] (-1.8, 0) -- (1.8, 0);
  \draw[->, thick] (0, -1.8) -- (0, 1.8);
  \draw[thick, engblue] (0,0) circle (1.5);
  \draw[->, very thick, engred] (0, 0) -- (1.5, 0) node[anchor=west, font=\scriptsize] {$\vect{v}_{1}$};
  \draw[->, very thick, engred] (0, 0) -- (0, 1.5) node[anchor=south, font=\scriptsize] {$\vect{v}_{2}$};
\end{scope}

% Arrow with V^T
\draw[->, thick] (1.9, 0) -- (2.6, 0) node[midway, above, font=\scriptsize] {$\mat{V}^{\transpose}$};

% Panel 2: rotated unit circle
\begin{scope}[shift={(4.2, 0)}]
  \node[anchor=south, font=\small\bfseries] at (0, 2.0) {rotate};
  \draw[->, thick] (-1.8, 0) -- (1.8, 0);
  \draw[->, thick] (0, -1.8) -- (0, 1.8);
  \draw[thick, engblue] (0,0) circle (1.5);
  \draw[->, very thick, engred] (0, 0) -- (1.5, 0) node[anchor=west, font=\scriptsize] {$\vect{e}_{1}$};
  \draw[->, very thick, engred] (0, 0) -- (0, 1.5) node[anchor=south, font=\scriptsize] {$\vect{e}_{2}$};
\end{scope}

\draw[->, thick] (6.1, 0) -- (6.8, 0) node[midway, above, font=\scriptsize] {$\mat{\Sigma}$};

% Panel 3: scaled ellipse, axes sigma_1 and sigma_2
\begin{scope}[shift={(8.6, 0)}]
  \node[anchor=south, font=\small\bfseries] at (0, 2.0) {scale};
  \draw[->, thick] (-2.4, 0) -- (2.4, 0);
  \draw[->, thick] (0, -1.4) -- (0, 1.4);
  \draw[thick, engblue] (0,0) ellipse (2.0 and 0.7);
  \draw[->, very thick, engred] (0, 0) -- (2.0, 0) node[anchor=west, font=\scriptsize] {$\sigma_{1}\vect{e}_{1}$};
  \draw[->, very thick, engred] (0, 0) -- (0, 0.7) node[anchor=south, font=\scriptsize] {$\sigma_{2}\vect{e}_{2}$};
\end{scope}

\draw[->, thick] (11.0, 0) -- (11.7, 0) node[midway, above, font=\scriptsize] {$\mat{U}$};

% Panel 4: final ellipse in output basis u_1, u_2
\begin{scope}[shift={(13.6, 0)}]
  \node[anchor=south, font=\small\bfseries] at (0, 2.0) {output ellipse};
  \draw[->, thick] (-2.4, 0) -- (2.4, 0);
  \draw[->, thick] (0, -1.6) -- (0, 1.6);
  \draw[thick, engblue, rotate=30] (0,0) ellipse (2.0 and 0.7);
  \draw[->, very thick, engred, rotate=30] (0, 0) -- (2.0, 0) node[anchor=west, font=\scriptsize] {$\sigma_{1}\vect{u}_{1}$};
  \draw[->, very thick, engred, rotate=30] (0, 0) -- (0, 0.7) node[anchor=south, font=\scriptsize] {$\sigma_{2}\vect{u}_{2}$};
\end{scope}

\node[font=\scriptsize, text width=14cm, align=center, anchor=north] at (6.5, -2.4) {
  $\mat{A} = \mat{U}\mat{\Sigma}\mat{V}^{\transpose}$ acts on the unit circle by
  rotating into the input singular basis, scaling each axis by $\sigma_{i}$, and
  rotating into the output singular basis. The condition number is the aspect
  ratio of the output ellipse: $\kappa(\mat{A}) = \sigma_{1}/\sigma_{n}$.
};

\end{tikzpicture}
\caption[The SVD as rotate, scale, rotate]{The singular value
decomposition $\mat{A} = \mat{U}\mat{\Sigma}\mat{V}^{\transpose}$
factors a linear map into three geometric operations: rotate the
input by $\mat{V}^{\transpose}$, scale axes by the singular values
$\sigma_{1} \geq \cdots \geq \sigma_{r} > 0$, and rotate the output
by $\mat{U}$. The unit ball maps to an ellipsoid whose principal
axes are $\sigma_{i}\vect{u}_{i}$. A small $\sigma_{i}$ flattens the
ellipsoid in the $\vect{u}_{i}$ direction; the input direction
$\vect{v}_{i}$ is the one the matrix nearly annihilates, and the
condition number $\kappa(\mat{A}) = \sigma_{1}/\sigma_{n}$ is the
aspect ratio of the output ellipsoid.}
\label{fig:vol02:ch09:svd-geometry}
\end{figure}
